\documentclass[conference]{IEEEtran}

% ---- packages ----
\usepackage[T1]{fontenc}
\usepackage[utf8]{inputenc}
\usepackage{cite}
\usepackage{amsmath,amssymb,amsfonts}
\usepackage{graphicx}
\usepackage{booktabs}
\usepackage{array}
\usepackage{textcomp}
\usepackage{xcolor}
\usepackage{url}
\usepackage[hidelinks]{hyperref}

% nicer column type for wrapping text in tables
\newcolumntype{L}[1]{>{\raggedright\arraybackslash}p{#1}}

\begin{document}

\title{CacheCollapse: Reproducing, Detecting, and Mitigating\\
Cache-Stampede-Driven Indirect Denial of Service in\\
Containerized Web Applications}

\author{
\IEEEauthorblockN{Manal Aamir, Tayyaba Arshad, Noor Afza, Batool Fatima, Aun Zaidi}
\IEEEauthorblockA{Department of Software Engineering\\
NED University of Engineering and Technology, Karachi, Pakistan\\
SE-401 Information Security (Complex Engineering Problem), Spring 2026\\
\{SE-23070, SE-23073, SE-23075, SE-23076, SE-23095\}}
}

\maketitle

\begin{abstract}
Containerized web applications almost universally place an in-memory cache between
a stateless application tier and a relational database to absorb read traffic. This
design is efficient under steady load, but it hides a structural availability
weakness: when a popular cache entry expires, a burst of otherwise legitimate
concurrent requests misses the cache simultaneously and is forwarded, unthrottled,
to the database. The resulting \emph{cache stampede} (also called dog-piling or the
thundering-herd problem) converts a modest amount of normal-looking traffic into a
flood of redundant database queries. Because the attacker never contacts the
database directly and never violates any application-level rule, the event is an
\emph{Indirect Distributed Denial of Service} (IDDoS) that volume-based network
defenses do not catch. This paper presents \textbf{CacheCollapse}, a fully
containerized testbed and defense framework that (i)~reproduces the cache-stampede
IDDoS deterministically on a nine-container Docker Compose stack of
industry-standard services, (ii)~detects it in real time using an explainable
multi-signal \emph{rule-based} detector and an unsupervised \emph{IsolationForest}
machine-learning detector trained only on normal traffic, and (iii)~mitigates it
with a layered defense combining \emph{cache locking} (single-flight regeneration
with stale-while-revalidate), \emph{application-layer rate limiting}, and \emph{edge
rate limiting} at the reverse proxy. Across a three-experiment protocol, an
undefended attack drove database load from ${\sim}15$ to a peak of ${\sim}230$
queries/s, raised average response time from ${\sim}130$~ms to ${\sim}700$~ms, and
produced up to ${\sim}180$ failed requests/s, while both detectors fired. With the
layered defense enabled, database load fell by roughly \textbf{25$\times$} (to
${\sim}8$ queries/s), failed requests dropped to zero, and customer-visible latency
on a demonstration storefront fell to ${\sim}80$~ms. An ablation shows that cache
locking alone protects the database but not the proxy connection path, and that
edge rate limiting is required to keep legitimate users fast during the attack.
\end{abstract}

\begin{IEEEkeywords}
Cache stampede, thundering herd, indirect denial of service, application-layer
DDoS, containerized security, rate limiting, cache locking, anomaly detection,
isolation forest, availability.
\end{IEEEkeywords}

\section*{Highlights}
\begin{itemize}
\item A reproducible nine-container Docker Compose testbed turns an ordinary
Nginx--Flask--Redis--PostgreSQL stack into a controllable cache-stampede IDDoS
laboratory; an undefended attack raised peak database load ${\sim}15\times$ and
produced up to ${\sim}180$ failed requests/s from legitimate-looking traffic.
\item A multi-signal rule-based detector and an unsupervised IsolationForest
detector trained only on baseline traffic both identified the attack; a side-by-side
comparison quantifies their trade-offs in training cost, false alarms, and
interpretability.
\item A layered mitigation of cache locking with stale-while-revalidate, application
rate limiting, and edge rate limiting reduced database query load under attack by
${\sim}25\times$ and restored failed requests to zero.
\item An ablation isolates a deployment-critical result: cache locking protects the
database but not the proxy, so \emph{edge} rate limiting is necessary to keep
legitimate users responsive during the flood.
\end{itemize}

\section{Introduction}
The availability of a web service is as much a security property as its
confidentiality or integrity. Denial-of-service attacks remain among the most common
and damaging threats to online systems, and the community has produced extensive
taxonomies and defenses for network- and transport-layer flooding
\cite{mirkovic2004,zargar2013}. A harder class operates at the application layer: the
adversary issues syntactically and semantically valid requests that are individually
indistinguishable from legitimate traffic, so purely volume- or signature-based
filters are ineffective \cite{praseed2019,xie2009}. This paper studies a particularly
subtle member of that class in the setting most modern software is deployed in---the
containerized, multi-service web application.

Virtually every read-heavy web application interposes an in-memory cache, such as
Redis, between its application tier and its database, so that repeated requests for
the same object are answered from memory. The entry is given a time-to-live (TTL),
after which it expires and must be regenerated from the database. If a popular
(``hot'') key expires while many requests for it are in flight, all of them miss the
cache at the same instant and are forwarded to the database simultaneously. A single
expiry thus produces not one query but $N$ redundant queries, where $N$ is the
instantaneous concurrency on that key. This is a \emph{cache stampede},
\emph{dog-pile}, or \emph{thundering herd}, documented as a performance hazard in
large-scale systems \cite{vattani2015,nishtala2013}.

What is less often treated is that the same mechanism is a \emph{security}
vulnerability. An adversary who understands a target's caching behavior can
deliberately align a burst of requests for a single hot object with that object's
expiry, repeatedly, so the database is driven into overload by traffic that never
leaves the normal request distribution of the application. The database---not the web
server---is the victim, and the web application is merely the conduit. We call this an
\emph{Indirect Distributed Denial of Service} (IDDoS): indirect because the attacker
never addresses the victim component directly, and distributed because the
amplification is produced by the fan-out of concurrency rather than by the attacker's
own request volume. Because the attack rides on legitimate application logic rather
than a software defect, it cannot be patched away; it must be \emph{detected} and
\emph{mitigated} as an operational property of the system.

\subsection{Problem Statement and the Complex Engineering Challenge}
Defending this vulnerability in a containerized environment is a genuine
complex-engineering problem because its difficulty arises from the \emph{interaction}
of well-understood components rather than from any single exotic one. Three tensions
must be resolved simultaneously:
\begin{itemize}
\item \textbf{Detection without a volumetric signature.} The attack's request rate at
the web tier can be indistinguishable from a legitimate spike or flash crowd
\cite{xie2009}. A detector keyed on request volume alone will either miss the attack
or raise constant false alarms on benign bursts.
\item \textbf{Mitigation without denying service to legitimate users.} A defense that
simply rejects traffic to protect the database can itself become the denial of
service. The mitigation must shed the \emph{redundant} load while continuing to serve
genuine users quickly.
\item \textbf{Faithful, safe, repeatable experimentation.} Studying an availability
attack requires driving a realistic system into failure repeatedly and measuring
per-component behavior, which demands an isolated, resettable, production-like
environment rather than a single-process simulation.
\end{itemize}
These requirements are coupled: an overly aggressive guardrail degrades legitimate
users, while a cache-only fix that ignores the connection path leaves users slow even
when the database is healthy (Section~\ref{sec:results}). Resolving the problem
requires an integrated architectural and algorithmic design, not merely a faster
database or a bigger cache.

\subsection{Proposed Solution}
We present \textbf{CacheCollapse}, a framework addressing all three dimensions
together: (i)~a \emph{reproducible containerized testbed} that reconstructs the
Nginx $\rightarrow$ Flask $\rightarrow$ Redis $\rightarrow$ PostgreSQL request path on
an isolated Docker Compose network and generates both normal traffic and a
controllable stampede attack; (ii)~a \emph{dual detection subsystem} pairing an
explainable multi-signal rule with an unsupervised IsolationForest model trained only
on baseline behavior; and (iii)~a \emph{layered mitigation} combining single-flight
cache locking with stale-while-revalidate, an application-layer per-client rate
limiter, and an edge rate limiter at the reverse proxy that can be toggled live.

\subsection{Contributions}
\begin{itemize}
\item A \textbf{containerized IDDoS testbed} that reproduces the cache-stampede attack
deterministically on nine interacting industry-standard services, with a control
plane that switches traffic, attack, and each defense on and off live.
\item A \textbf{multi-signal rule-based detector} that fires only when elevated
database load, elevated latency, and an elevated cache-miss rate coincide, with
cold-cache warm-up suppression that eliminates start-up false alarms.
\item An \textbf{unsupervised IsolationForest detector} trained exclusively on
normal-baseline traffic with standardized features, and a controlled comparison of
the two detectors on training cost, responsiveness, false alarms, and
interpretability.
\item A \textbf{layered mitigation} and an \textbf{ablation} isolating each
mechanism's contribution, including the deployment-critical finding that
database-level and proxy-level protection are distinct and both required.
\item A fully \textbf{reproducible, instrumented deployment} (single
\texttt{docker compose up}, Prometheus/Grafana, live dashboard, demonstration
storefront) validated by a 16-case automated test suite.
\end{itemize}

\section{Literature Review}
\subsection{Cache Stampede and the Thundering Herd}
Nishtala et al., describing Facebook's memcache deployment, document ``stale sets''
and ``thundering herds'' as first-order concerns and introduce \emph{leases}---per-key
tokens that permit only one client to regenerate a missing value while others briefly
wait---reporting a large reduction in peak database query rate \cite{nishtala2013}.
Vattani, Chierichetti, and Lowenstein formalize the mitigation side, deriving
theoretically optimal \emph{probabilistic early recomputation} (XFetch), in which an
item is refreshed slightly before its TTL with a probability that rises as expiry
approaches \cite{vattani2015}. Both treat the stampede as a performance hazard.
CacheCollapse builds on their mitigation insights---single-flight regeneration is the
security-hardened descendant of memcache leases---but reframes the stampede as an
\emph{attacker-triggerable} availability vulnerability and adds detection and
operational defense.

\subsection{Denial of Service and Application-Layer Attacks}
Classical DDoS research classifies attacks and defenses mainly at the network and
transport layers. Mirkovic and Reiher give the canonical taxonomy \cite{mirkovic2004},
and Zargar et al. survey flooding defenses, stressing that protection must span
prevention, detection, and response \cite{zargar2013}. Application-layer DoS is a
distinct, harder problem: Xie and Yu model normal browsing to detect app-layer attacks
using ordinary HTTP GET requests \cite{xie2009}, and Praseed and Thilagam survey the
space, noting such attacks can be mounted at low volume with legitimate requests
\cite{praseed2019}. The cache-stampede IDDoS is an extreme point: not only are the
requests legitimate, the \emph{aggregate} web-tier rate need not be anomalous---the
damage is produced downstream by cache semantics.

\subsection{Availability Attacks in Containerized Environments}
Containers are the dominant packaging and orchestration unit for modern applications
\cite{pahl2015,merkel2014}, with cluster managers such as Kubernetes tracing their
lineage to internal systems like Borg and Omega \cite{burns2016}, and the microservice
style decomposing applications into interacting services \cite{dragoni2017}. Security
analyses of Docker catalog isolation and configuration weaknesses
\cite{combe2016,sultan2019}, but availability has received comparatively less
attention. A directly relevant example is the \emph{Yo-Yo attack}, in which periodic
bursts drive an auto-scaling controller to oscillate, degrading performance and
inflating cost while evading simple volume detection \cite{sides2016}. CacheCollapse
is complementary: like Yo-Yo it exploits an emergent property of a cloud-native
pattern (there, elasticity; here, caching) using legitimate-looking traffic, and must
be studied in an orchestrated, resettable environment.

\subsection{Detection: Rule-Based and Machine-Learning Anomaly Detection}
Anomaly detection is the natural framing for an attack with no fixed signature.
Chandola et al. organize the field by technique and assumption, underlining that
detector choice depends on label availability and interpretability requirements
\cite{chandola2009}. Among unsupervised methods, the IsolationForest of Liu, Ting, and
Zhou isolates anomalies directly with randomized trees, has near-linear complexity,
needs no labels, and performs well with small samples \cite{liu2008}; we use its
scikit-learn implementation \cite{pedregosa2011}. CacheCollapse deliberately pairs an
unsupervised model with an explainable rule because the two occupy different points
in the design space \cite{chandola2009}.

\subsection{Mitigation: Rate Limiting and Request Coalescing}
Raghavan et al. study distributed rate limiting that enforces a global aggregate limit
while preserving fair flow behavior \cite{raghavan2007}, and production proxies expose
mature local limiters---Nginx's \texttt{limit\_req} implements a leaky-bucket limiter
keyed on client address \cite{nginxlimitreq}. Request coalescing is realized through
Redis TTLs and \texttt{SET\ldots NX} primitives \cite{redisdocs} and is conceptually
descended from memcache leases \cite{nishtala2013}; stability-pattern practice further
recommends serving stale data rather than blocking during regeneration
\cite{nygard2018}. A contribution of this paper is to show empirically that these
controls are \emph{not} interchangeable: coalescing protects the database, rate
limiting protects the request path, and only their combination preserves end-to-end
availability (Section~\ref{sec:ablation}).

\section{Methodology}
\subsection{Threat Model and Attack Chain}
We assume an external adversary who can send ordinary HTTP requests to the public
endpoint and knows (or can probe) that responses for a set of hot objects are cached
with a short TTL. The adversary cannot reach the database or cache directly, cannot
read internal metrics, and issues only well-formed requests. The attack chain is:
\emph{hot entry expires} $\rightarrow$ \emph{$N$ concurrent requests arrive}
$\rightarrow$ \emph{$N$ cache misses} $\rightarrow$ \emph{$N$ database queries}
$\rightarrow$ \emph{overload}.

Let a hot key $k$ have regeneration time $d$. If $N$ requests for $k$ arrive within
$(t_{\mathrm{exp}}, t_{\mathrm{exp}}+d)$, an unprotected system issues $N$ database
queries; an ideal single-flight system issues $1$. The ratio $N{:}1$ is the stampede
amplification, which the attacker maximizes by concentrating requests on hot keys and
phasing bursts to expiry. Our testbed uses a short TTL and an attacker that rotates
its burst across a small pool of hot keys, so some key is always expiring.

\subsection{System Architecture}
CacheCollapse is a nine-service Docker Compose application on one isolated bridge
network; each service is a separate container so CPU, memory, and I/O can be
attributed per component and the system reset between runs (Table~\ref{tab:services}).
The request path is Nginx $\rightarrow$ Flask $\rightarrow$ Redis, with a PostgreSQL
read only on a cache miss. Redis additionally serves as a lightweight \emph{control
bus}: the dashboard writes run-state and defense flags that the load generator and
Flask app read, so every experiment is reproducible from the browser.

\begin{table}[t]
\centering
\caption{CacheCollapse containerized services.}
\label{tab:services}
\footnotesize
\begin{tabular}{@{}cL{1.3cm}L{4.8cm}@{}}
\toprule
\# & Service & Role \\
\midrule
1 & Nginx & Public reverse proxy; edge (\texttt{limit\_req}) rate limiting \\
2 & Flask app & \texttt{GET /products/<id>}; cache locking + app rate limiting \\
3 & Redis & Cache (short TTL) and shared control bus \\
4 & PostgreSQL & Victim database; 100 seeded product records \\
5 & Load gen. & Asynchronous normal traffic and the burst attack \\
6 & Monitor & Per-second metrics, rule + ML detection, CSV log, live feed \\
7 & Frontend & Live dashboard and demonstration storefront \\
8 & Prometheus & Time-series metrics collection \\
9 & Grafana & Standard dashboards over Prometheus \\
\bottomrule
\end{tabular}
\end{table}

\subsection{Attack and Normal-Traffic Generation}
The load generator reads two flags from the control bus. With \emph{normal traffic}
on, it issues a steady, low-rate stream spread across the catalog, representing a
benign population whose requests mostly hit the cache. With \emph{attack} on, it
issues large concurrent bursts aimed at a small pool of hot keys, rotating the target
between bursts so the pool continuously cycles through expiry. This rotation converts
an isolated, one-shot stampede into a \emph{sustained} overload.

\subsection{Instrumentation and Metrics}
The monitor samples twice per second: it scrapes the Flask app's Prometheus metrics,
computes per-interval rates, runs both detectors, appends a CSV row, and streams the
result to the dashboard over Server-Sent Events. Measured quantities are requests/s,
database queries/s, cache hits/misses (and the miss rate $m=\text{misses}/(\text{hits}
+\text{misses})$), average and p95 response time, failed requests/s, and per-container
CPU and memory. Two timings are recorded per attack: \emph{detect time} (onset to
first alert) and \emph{recover time} (defense activation to healthy state). Failures
are split into \texttt{HTTP 503} (genuine overload) and \texttt{HTTP 429}
(intentional shedding by the rate limiter); only the former count as availability
failures, so an active defense is not mis-scored as an outage.

\subsection{Rule-Based Detector}
The rule encodes the attack's multi-signal signature, firing only when elevated
database load, latency, and cache-miss rate hold simultaneously:
\begin{equation}
\mathrm{alert} \iff
(q_{db} > K\,\bar{q}_{\mathrm{base}}) \wedge
(\bar{L} > \tau_L) \wedge
(m > \tau_m),
\label{eq:rule}
\end{equation}
where $q_{db}$ is database queries/s, $\bar{q}_{\mathrm{base}}$ a learned baseline,
$\bar{L}$ average response time, and $m$ the cache-miss rate; we use $K=5$,
$\tau_L=200$~ms, $\tau_m=0.5$. Requiring all three conditions distinguishes an attack
from a benign flash crowd: a legitimate spike raises request volume but is served
from cache, so it does not simultaneously raise database load and the miss rate. A
cold-cache warm-up suppression disables alerting during the first few seconds after
traffic starts, and $\bar{q}_{\mathrm{base}}$ is learned online from steady traffic.

\subsection{Machine-Learning Detector}
The ML detector is an unsupervised IsolationForest \cite{liu2008,pedregosa2011} trained
only on normal traffic---no attack samples and no labels. Each sample is
$\mathbf{x}=[\,q_{db}, \bar{L}, m, \text{req/s}, \text{failed/s}\,]$. Because features
span different scales, they are standardized to zero mean and unit variance using
baseline statistics before the forest is fit; live samples are standardized with the
same statistics and scored, and a sample flagged as an outlier by the forest's
decision function is an anomaly. The model trains only after a short warm-up, and an
alert must persist for two consecutive samples (debounce), removing single-interval
blips while leaving a sustained attack easily detectable.

\subsection{Mitigation I---Cache Locking with Stale-While-Revalidate}
On a miss for key $k$, the Flask app attempts a short-lived mutex via
\texttt{SET lock:}$k$\texttt{ 1 NX PX T}. The single winner queries the database and
writes a \emph{fresh} copy under the normal TTL and a \emph{stale} copy under a much
longer TTL. Every other concurrent miss immediately returns the stale copy. Thus $N$
simultaneous misses produce exactly \emph{one} database query, and losers are served
in cache time rather than waiting for regeneration. A stale copy is pre-warmed for
every key at start-up so the defense is instant even if an attack begins before any
value is cached---the security-oriented generalization of memcache leases
\cite{nishtala2013} and of stale-serving stability practice \cite{nygard2018}.

\subsection{Mitigation II---Layered Rate Limiting}
Rate limiting sheds the attacker's excess concurrency at two layers. The
\emph{application layer} enforces a per-client, per-second cap inside Flask using
Redis counters, returning \texttt{HTTP 429} to excess requests; running behind the
proxy, it is toggled live without a reload. The \emph{edge layer} uses Nginx's
leaky-bucket \texttt{limit\_req} keyed on client address \cite{nginxlimitreq}; when
enabled, the monitor rewrites a small included configuration fragment and triggers an
Nginx reload, so the flood is rejected at the proxy before it consumes application or
connection resources. The proxy forwards the real client address to the application
so the per-client limiter buckets the attacker and a genuine user separately.

\subsection{Experimental Protocol}
Experiments follow a three-condition protocol executed from the dashboard:
\textbf{E1---Baseline} (normal traffic only); \textbf{E2---Attack, undefended};
and \textbf{E3---Attack, defended} (cache locking, application rate limiting, and edge
rate limiting enabled). Each condition is sampled for a fixed window, and all metrics
are logged. A demonstration storefront (``CacheShop'') that loads hot products and
reports per-item load time provides a customer's-eye view of availability. All
experiments ran on a single host under Docker Desktop (Docker~29.x).

\section{Results and Analysis}
\label{sec:results}
\subsection{Setup}
The database was seeded with 100 records; a short TTL and a rotating-hot-key attacker
were used so stampedes recur and overload is sustained. Reported figures are
representative values from the logs; peaks and interval averages are distinguished
where relevant.

\subsection{Cache Behavior (Sanity Check)}
The first request for a product is answered from the database in ${\sim}180$~ms, and
the immediately following request for the same product is answered from Redis in
${\sim}12$~ms---a ${\sim}15\times$ speed-up confirming the cache is the component
under contention during a stampede.

\subsection{Three-Experiment Comparison}
Table~\ref{tab:three} summarizes the conditions. Under the undefended attack (E2),
database load rose from ${\sim}15$ to a peak of ${\sim}230$ queries/s, average
response time from ${\sim}130$~ms to ${\sim}700$~ms, the miss rate toward ${\sim}0.88$,
and up to ${\sim}180$ requests/s failed with \texttt{HTTP 503}---a clear overload from
legitimate-looking requests. With the layered defense (E3), database load fell to
${\sim}8$ queries/s (a ${\sim}25\times$ reduction), response time to ${\sim}90$~ms, and
\texttt{503} failures to zero.

\begin{table}[t]
\centering
\caption{Three-experiment comparison (representative values).}
\label{tab:three}
\footnotesize
\begin{tabular}{@{}L{2.0cm}cccc@{}}
\toprule
Condition & DB q/s (pk) & Resp. (pk) & Fail/s & Status \\
\midrule
E1 Baseline & ${\sim}15$ & ${\sim}130$~ms & 0 & Healthy \\
E2 Undefended & ${\sim}230$ & ${\sim}700$~ms & ${\sim}180$ & Overloaded \\
E3 Defended & ${\sim}8$ & ${\sim}90$~ms & 0 & Protected \\
\bottomrule
\end{tabular}
\end{table}

\subsection{Detection: Rule-Based vs. Machine-Learning}
Both detectors identified the attack in E2 and stayed silent in E1
(Table~\ref{tab:detectors}). The rule requires no training, is immediately available,
is fully interpretable (an alert names which conditions held), and did not fire on
calm traffic; but because the attack is bursty, the three conditions align only
intermittently, so the rule fires on a subset of attack intervals. The IsolationForest,
once trained on a short baseline, fired on essentially every attack interval,
providing denser, faster confirmation, at the cost of requiring a clean training
window and offering less transparent justification. Used together, the rule explains
\emph{why} an alert occurred while the forest provides continuous confirmation.

\begin{table}[t]
\centering
\caption{Detector comparison.}
\label{tab:detectors}
\footnotesize
\begin{tabular}{@{}L{2.5cm}L{2.3cm}L{2.3cm}@{}}
\toprule
Property & Rule-based & IsolationForest \\
\midrule
Training data & None & Baseline only \\
Availability & Immediate & After warm-up \\
Baseline false alarms & None & Rare (debounced) \\
Attack-interval coverage & Intermittent & Dense \\
Interpretability & High & Low \\
\bottomrule
\end{tabular}
\end{table}

\subsection{Mitigation Ablation}
\label{sec:ablation}
An ablation isolates each defense (Table~\ref{tab:ablation}) and yields the paper's
central deployment finding. Cache locking alone \emph{protects the database}---query
load drops to single digits because $N$ misses are coalesced to one---but does not by
itself keep users fast: with only application-layer controls, the attacker's flood
still congests the reverse proxy and connection path, so the storefront remained slow
even though the database was healthy. Stale-while-revalidate removed the regeneration
queue; edge rate limiting removed the remaining congestion by rejecting the flood at
the proxy, and only then did storefront latency fall to ${\sim}80$~ms. Database
protection and request-path protection are thus distinct problems requiring distinct
controls.

\begin{table}[t]
\centering
\caption{Mitigation ablation (customer-visible behavior under attack).}
\label{tab:ablation}
\footnotesize
\begin{tabular}{@{}L{2.4cm}L{1.7cm}L{1.6cm}L{1.4cm}@{}}
\toprule
Configuration & DB load & Storefront & Avail. \\
\midrule
No defense & Overloaded & Fails/slow & Denied \\
Cache lock only & Protected & Still slow & Degraded \\
+ Stale-while-reval. & Protected & Improved & Degraded \\
+ Edge rate limit & Protected & ${\sim}80$~ms & Preserved \\
\bottomrule
\end{tabular}
\end{table}

\subsection{Recovery and Rate-Limiting Layers}
When defenses are enabled during an ongoing attack, the system returns to a healthy
state within ${\sim}1.5$--$2.5$~s, because cache locking immediately collapses the
redundant queries and the pre-warmed stale cache serves losers without waiting.
Isolating the rate-limiting layers confirms they protect different things:
application-layer limiting sheds excess \emph{database} work (returning \texttt{429}
rather than touching the database) but does not relieve proxy congestion, whereas edge
\texttt{limit\_req} rejects the flood at the proxy and brings storefront latency down
by roughly an order of magnitude---the empirical basis for recommending edge rate
limiting, not application limiting alone.

\subsection{End-to-End Validation}
A 16-case automated test suite exercised against the live stack checks service health
and routing, cache behavior, the calm baseline (no false alarms), the undefended
attack (overload with both detectors firing), each mitigation, recovery after the
attack stops, and the live-feed rate. All 16 cases pass, providing repeatable evidence
for the results above.

\section{Comparative Analysis}
CacheCollapse integrates concerns that prior work addresses only in part
(Table~\ref{tab:capability}). Performance-oriented stampede work
\cite{vattani2015,nishtala2013} contributes coalescing but no attacker model,
detection, or testbed. Application-layer DDoS work \cite{praseed2019,xie2009}
contributes detection but neither the cache-specific mechanism nor a reproducible
container environment. Container-security work \cite{combe2016,sultan2019} catalogs
weaknesses without instrumenting a specific availability attack end to end. The Yo-Yo
attack \cite{sides2016} is closest in spirit but targets elasticity rather than caching
and does not combine detection with layered mitigation. CacheCollapse is the only one
of these to address all four dimensions together.

\begin{table}[t]
\centering
\caption{Capability coverage. \checkmark~= addressed; $\sim$~= partial; $\times$~= not addressed.}
\label{tab:capability}
\footnotesize
\begin{tabular}{@{}L{2.9cm}cccc@{}}
\toprule
Work & Testbed & Detect & Layered mit. & Containerized \\
\midrule
Prob. prevention \cite{vattani2015} & $\times$ & $\times$ & $\sim$ & $\times$ \\
Memcache leases \cite{nishtala2013} & $\times$ & $\times$ & $\sim$ & $\times$ \\
App-layer DDoS \cite{praseed2019,xie2009} & $\times$ & \checkmark & $\times$ & $\times$ \\
DDoS surveys \cite{mirkovic2004,zargar2013} & $\times$ & $\sim$ & $\sim$ & $\times$ \\
Container security \cite{combe2016,sultan2019} & $\times$ & $\times$ & $\times$ & $\sim$ \\
Yo-Yo attack \cite{sides2016} & $\sim$ & $\sim$ & $\times$ & \checkmark \\
\textbf{CacheCollapse} & \checkmark & \checkmark & \checkmark & \checkmark \\
\bottomrule
\end{tabular}
\end{table}

\section{Discussion}
\subsection{Why a Layered Defense Is Necessary}
No single control is sufficient. Cache locking is the only mechanism that addresses the
root cause---it restores the $N{:}1$ query ratio to $1{:}1$---but operates below the
proxy and cannot relieve edge congestion. Rate limiting relieves that congestion but,
applied only at the application layer, acts too late in the path; applied at the edge
it protects the proxy but does not fix the redundant-query amplification if a limiter
is tuned permissively. The correct design composes them: edge limiting caps the flood,
cache locking coalesces the surviving misses, and stale-while-revalidate serves the
coalesced waiters without delay---mirroring the prevention--detection--response framing
urged for DDoS defense \cite{zargar2013}.

\subsection{Rules vs. Machine Learning}
The comparison argues against treating the detectors as substitutes. The rule is the
better \emph{first} line (immediate, interpretable, free of training and false-alarm
risk on benign spikes); the forest is the better \emph{continuous} confirmation (dense
firing, adaptive to baseline). Their weaknesses are complementary, which is why
deploying both is preferable to choosing one \cite{chandola2009}.

\subsection{Deployment Implications}
For operators of cached, containerized web services: enable single-flight regeneration
with a stale-serving fallback on hot keys; rate-limit at the edge, not only in the
application; forward real client identity so per-client limits are meaningful; and
monitor database load, latency, and cache-miss rate \emph{together}, since any one
alone is an unreliable alarm. All are realizable with standard components---Redis,
Nginx, and a reverse proxy.

\subsection{Limitations}
The testbed runs on a single host under Docker Desktop, whose port-forwarding adds some
latency under heavy connection load, so absolute latency figures are
environment-dependent while the \emph{relative} effects are robust. The normal
population and attacker are synthetic; validating against production traces would
strengthen the evaluation. The IsolationForest requires a clean baseline window, so a
system attacked from start-up would train on contaminated data; the always-available
rule mitigates but does not eliminate this. Finally, thresholds ($K$, $\tau_L$,
$\tau_m$, limiter rates) were tuned on the testbed and need site-specific calibration.

\section{Conclusion}
CacheCollapse is a reproducible containerized framework for studying and defending the
cache-stampede IDDoS---an availability attack in which legitimate-looking requests,
aligned with cache expiry, overload a database through the application rather than
attacking it directly. On a nine-container stack, an undefended attack raised peak
database load roughly $15\times$, pushed average response time from ${\sim}130$~ms to
${\sim}700$~ms, and produced up to ${\sim}180$ failed requests/s. A multi-signal rule
and an unsupervised IsolationForest each identified the attack while remaining silent
on benign traffic, and their comparison clarified a real trade-off between
interpretability and continuous sensitivity. A layered mitigation reduced database
query load by roughly $25\times$, eliminated \texttt{503} failures, and restored fast
response, recovering within a couple of seconds; an ablation established that
database-level coalescing and edge-level rate limiting protect different parts of the
system and are both required. Future work will replay production traces, extend the
testbed to multi-node Kubernetes to study interaction with auto-scaling attacks such
as Yo-Yo \cite{sides2016}, add probabilistic early recomputation \cite{vattani2015} as
a proactive complement, and automate threshold calibration.

\section*{CRediT Authorship Contribution Statement}
\textbf{Manal Aamir:} Conceptualization, Methodology, Writing -- Review and Editing.
\textbf{Tayyaba Arshad:} Investigation, Validation, Data Curation.
\textbf{Noor Afza:} Software (load generation and monitoring), Visualization.
\textbf{Batool Fatima:} Software (detection and mitigation), Validation.
\textbf{Aun Zaidi:} Conceptualization, Methodology, Software (system integration),
Writing -- Original Draft.

\bibliographystyle{IEEEtran}
\bibliography{references}

\end{document}
